\chapter[1937 --- Szent-Gyorgyi]{1937: Albert Szent-Gy\"orgyi}
\label{ch:1937_albertszentgyorgyi}

\section*{Main Contribution}

\textit{Identified vitamin C (ascorbic acid) and demonstrated its role in preventing scurvy and in collagen synthesis, advancing understanding of nutritional biochemistry.}

\section{Official Citation}

for his discoveries in connection with the biological combustion of vitamins

\section{Background and Historical Context}

Albert Szent-Gy\"orgyi was born on September 16, 1893, in Budapest, Hungary, into a distinguished family with a long tradition of intellectual achievement. His father, Miklós Szent-Gy\"orgyi, was a prominent lawyer and landowner who also served as a member of the Hungarian Parliament. His mother, Josephine von Lósztó, came from a wealthy family, and young Albert grew up in an environment that valued education and scientific inquiry. From an early age, Szent-Gy\"orgyi showed a remarkable curiosity about the natural world, spending hours examining plants and insects in the gardens of the family estate.

Szent-Gy\"orgyi's formal education began in Budapest, where he attended the Piarist Gymnasium before enrolling at the University of Budapest in 1911. His studies were interrupted by World War I, during which he served as a medical officer in the Austro-Hungarian Army. The horrors of war left a significant impression on him, and after the conflict, he returned to his studies with renewed determination. He transferred to the University of Szeged in 1920, where he would spend much of his early career. Under the mentorship of several prominent scientists, Szent-Gy\"orgyi developed his interest in biochemistry and began investigating the chemical processes that occur within living organisms.

The scientific landscape of the early twentieth century was undergoing rapid transformation. The discovery of vitamins was still in its infancy, and researchers were beginning to understand that certain dietary factors were essential for human health. Scurvy, a disease that had plagued sailors for centuries, was known to be linked to citrus fruits, but the exact chemical compound responsible for preventing the disease remained unidentified. Carl and Gerty Cori, working in the United States, had recently elucidated the glycogen cycle, while other researchers were making progress in understanding the role of organic acids in cellular metabolism. It was within this vibrant scientific environment that Szent-Gy\"orgyi made his pioneering discoveries.

\section{The Discovery/Contribution}

Szent-Gy\"orgyi's Nobel Prize was awarded ``for his discoveries in connection with the biological combustion processes, with special reference to vitamin C and the catalysis of fumaric acid.'' His work spanned two major areas of investigation that were, as he would demonstrate, intimately connected: the role of ascorbic acid (vitamin C) in biological systems and the importance of fumaric acid and other dicarboxylic acids in cellular respiration.

\subsection{The Isolation of Ascorbic Acid}

Szent-Gy\"orgyi's journey toward the isolation of vitamin C began in the early 1920s, when he was working at the University of Szeged. His initial research focused on biological oxidation, specifically the process by which cells convert nutrients into energy. While investigating the respiratory pigments of adrenal glands, he observed that the adrenal cortex contained a remarkably high concentration of a reducing substance---a compound capable of donating electrons in chemical reactions. Through careful experimentation, he demonstrated that this reducing substance was distinct from known compounds and that it was present in many plant and animal tissues.

In 1927, while working at the Institute of Chemistry in Cambridge, England, Szent-Gy\"orgyi isolated a hexuronic acid from adrenal gland tissue. This compound, which he initially called ``hexuronic acid,'' was found to have potent reducing properties and appeared to be widely distributed in biological materials. Working with Joseph Svirbely, he later demonstrated that this same compound was identical to the antiscorbutic factor---the substance responsible for preventing scurvy. In 1932, Szent-Gy\"orgyi and Svirbely published their landmark paper identifying hexuronic acid as vitamin C, which they renamed ``ascorbic acid'' (from ``a'' meaning ``without'' and ``scorbutus'' meaning ``scurvy'').

The chemical structure of ascorbic acid was determined through the combined efforts of Szent-Gy\"orgyi's group and that of Norman Haworth at the University of Birmingham. In 1933, Szent-Gy\"orgyi achieved the first total synthesis of vitamin C, a feat that would later be exploited for industrial production. His synthesis was remarkably elegant, beginning with the natural sugar L-gulose and involving a series of carefully controlled chemical transformations. The ability to synthesize vitamin C was of enormous practical significance, as it meant that the vitamin could be produced in large quantities without relying on extraction from natural sources.

\subsection{The Role of Fumaric Acid in Cellular Respiration}

Szent-Gy\"orgyi's other major contribution concerned the role of dicarboxylic acids in cellular respiration---the process by which cells generate energy from nutrients. Working with slices of pigeon breast muscle, he observed that small amounts of certain four-carbon dicarboxylic acids, particularly fumaric acid, oxaloacetic acid, and succinic acid, could stimulate oxygen consumption far beyond what would be expected from the quantity of substrate added. This ``catalytic'' effect was puzzling and led Szent-Gy\"orgyi to propose that these acids were not merely substrates being consumed in oxidation but were instead participants in a catalytic cycle.

Szent-Gy\"orgyi's insight was that fumaric acid and its related compounds formed a cycle in which they were sequentially transformed by enzymatic reactions. In this cycle, oxaloacetic acid was reduced to malic acid, which was then dehydrated to fumaric acid. Fumaric acid was in turn reduced to succinic acid, which was then oxidized back to oxaloacetic acid, completing the cycle. This cycle, which Szent-Gy\"orgyi described in a series of papers published between 1935 and 1937, represented a significant departure from the prevailing view that cellular respiration was a linear process.

Szent-Gy\"orgyi's proposed cycle was not entirely complete---the steps connecting succinic acid to oxaloacetic acid would later be filled in by Krebs and others---but it contained the essential core of what would become known as the citric acid cycle (or Krebs cycle). His demonstration that small amounts of dicarboxylic acids could catalyze the oxidation of large quantities of carbohydrate was a fundamental insight that opened up an entirely new field of biochemistry. The Nobel Committee recognized the significant importance of these discoveries by awarding Szent-Gy\"orgyi the Prize in 1937.

\section{Impact on Modern Medicine}

Szent-Gy\"orgyi's discoveries have had far-reaching consequences for modern medicine and biology. The identification and synthesis of vitamin C revolutionized our understanding of nutrition and disease prevention. Vitamin C is now known to be essential for the synthesis of collagen, the most abundant protein in the human body, which provides structural support to skin, blood vessels, bones, and connective tissues. Without adequate vitamin collagen production is impaired, leading to the characteristic symptoms of scurvy: bleeding gums, poor wound healing, and ultimately death.

The importance of vitamin C extends well beyond the prevention of scurvy. As a powerful antioxidant, vitamin C protects cells from damage caused by free radicals---highly reactive molecules that are generated during normal metabolic processes and in response to environmental stressors such as pollution, ultraviolet radiation, and cigarette smoke. By neutralizing free radicals, vitamin C helps prevent the oxidative damage that is implicated in a wide range of diseases, including cardiovascular disease, cancer, and neurodegenerative disorders.

Vitamin C also plays a crucial role in immune function. It stimulates the production and activity of white blood cells, enhances the function of the immune barrier in the skin and mucous membranes, and supports the production of antibodies. During the COVID-19 pandemic, vitamin C supplementation was widely studied as a potential adjunctive therapy, and while the evidence for its efficacy in treating acute infection remains mixed, its role in supporting immune function is well established.

Perhaps most significantly for modern medicine, Szent-Gy\"orgyi's work on fumaric acid and the citric acid cycle laid the foundation for our understanding of cellular metabolism. The citric acid cycle is the central metabolic pathway in aerobic organisms, responsible for the oxidation of acetyl-CoA derived from carbohydrates, fats, and proteins. The energy released during this process is captured in the form of reduced coenzymes (NADH and FADH2), which subsequently donate electrons to the electron transport chain, driving the synthesis of ATP---the universal energy currency of the cell.

Understanding the citric acid cycle has been essential for the development of modern medicine. Inborn errors of metabolism affecting enzymes of the cycle cause a range of genetic diseases, and understanding these disorders has led to improved diagnostic and therapeutic approaches. The citric acid cycle also plays a central role in cancer metabolism: the Warburg effect, in which cancer cells preferentially use glycolysis even in the presence of oxygen, involves alterations in the citric acid cycle and its associated pathways. Current research into cancer metabolism is drawing heavily on Szent-Gy\"orgyi's foundational discoveries.

Szent-Gy\"orgyi himself went on to have a remarkably productive career. He continued his research on muscle contraction and discovered actin and myosin, the proteins responsible for muscle movement. He also became a passionate advocate for science and education, founding the Institute for Muscle Research at the Woods Hole Marine Biological Laboratory in Massachusetts. He was elected to the National Academy of Sciences and received numerous honors during his long career, which extended until his death in 1986 at the age of 93.

\section{Legacy}

Albert Szent-Gy\"orgyi's legacy is multifaceted and enduring. His identification and synthesis of vitamin C transformed our understanding of nutrition and made possible the mass production of this essential nutrient. His work on fumaric acid and the citric acid cycle opened up an entirely new field of biochemistry and laid the groundwork for our modern understanding of cellular metabolism. Beyond his scientific achievements, Szent-Gy\"orgyi was a forward-thinking researcher who saw the connections between disparate areas of research and who was not afraid to challenge prevailing orthodoxies.

The institutions that Szent-Gy\"orgyi founded and the students he mentored continued to advance the fields he helped create. The Institute for Muscle Research at Woods Hole became a major center for the study of muscle biology, and many of Szent-Gy\"orgyi's former students went on to distinguished careers in biochemistry and medicine. His published works, including several influential textbooks, have educated generations of scientists and remain valuable references to this day.

Szent-Gy\"orgyi's story also serves as a reminder of the importance of basic research. His discoveries were not made in pursuit of any specific medical application; they were driven by a fundamental curiosity about the workings of living systems. Yet these discoveries have had significant and practical consequences for human health, demonstrating the unpredictable but enormous value of investing in basic science. As Szent-Gy\"orgyi himself once remarked, ``Discovery consists of seeing what everybody has seen and thinking what nobody has thought.'' This philosophy guided his work throughout his career and continues to inspire scientists around the world.

Today, as we face new challenges in medicine---from the emergence of novel infectious diseases to the growing burden of chronic degenerative conditions---the principles underlying Szent-Gy\"orgyi's discoveries remain as relevant as ever. Our understanding of metabolism, nutrition, and oxidative stress continues to evolve, building on the foundation he laid nearly a century ago. The 1937 Nobel Prize in Physiology or Medicine thus represents not merely a recognition of past achievements but a milestone in an ongoing journey of discovery that continues to improve human health and well-being.


\section{Modern Developments}

Szent-Gyorgyi's vitamin C work has evolved into modern nutritional biochemistry. Understanding of collagen synthesis has led to treatments for scurvy and wound healing. High-dose vitamin C is being studied as an adjunct to chemotherapy, with mixed results. The role of vitamin C in immune function has been highlighted during the COVID-19 pandemic. Liposomal vitamin C formulations improve bioavailability. Understanding of oxidative stress has led to development of targeted antioxidant therapies.


\section{Bibliography}

\begin{thebibliography}{9}
\bibitem{nobel-1937}
Nobel Prize Outreach, ``The Nobel Prize in Physiology or Medicine 1937,''\emph{NobelPrize.org}.\\
\url{https://www.nobelprize.org/prizes/medicine/1937/summary/}

\bibitem{lecture-1937}
Albert Szent-Gy"orgyi, ``Nobel Lecture,''\emph{NobelPrize.org}. Winner-authored primary source; downloadable from the official Nobel Prize archive.\\
\url{https://www.nobelprize.org/prizes/medicine/1937/szent-gyorgyi/lecture/}
\end{thebibliography}
